\subsection{Historický vývoj systému}
\label{subsec:kelvin-history}

Systém Kelvin vznikl původně jako obyčejný školní projekt od studenta Daniel Trnka \footnote[1]{https://github.com/trnila}, ale v průběhu let byl postupně rozšiřován dalšími studenty a vyučujícími.
Tento způsob vývoje bez jednotné dlouhodobé architektonické vize vedl ke vzniku technického dluhu.
Ten se v systému projevuje například ve formě dlouhých a obtížně čitelných souborů, nedostatečného nebo chybějícího typování a míchání různých programátorských přístupů.

Uvedené skutečnosti mají výrazný vliv na implementaci nových funkcionalit.
Před přidáním nové aplikační logiky je často nutné provést částečný refaktoring existujícího kódu, aby byla zachována alespoň základní úroveň čitelnosti a udržovatelnosti systému.
Tato nutnost dále zpomaluje vývoj, ale zároveň je nezbytná pro dlouhodobou stabilitu a rozšiřitelnost systému Kelvin.

\endinput